\section{Preparation of the two-endpoint statuses}
\label{sec:preparation}

We work along a sequence of unit laws with endpoint marginals at
most \(M\pi\) and projected pair density at most
\(2^{3000gN}\pi^2\), as in the first part of
\Cref{thm:distribution-tests}.  A colour is an additional mark and
may be randomized conditional on the unit.  Passing to subsequences
is allowed.  By \Cref{lem:channel-thinning}, the absolute joint
density relative to \(\rawlaw^2\) has ample slack below
\(2^{DN}\).  Apply \Cref{cor:coloured-peeling}, or
\Cref{lem:peeling} when there is no colour, and discard the
exception in \Cref{lem:raw-intersections}.  This removes only
\(o(1)\) total mass.

All restrictions in this section act on the resulting \emph{pooled
law}, with the actual leaf weights.  Although a colour is fixed on
each leaf, no prediction or phase argument is applied separately to
its normalized colour class.  In particular the endpoint marginal
bound depends on the total retained mass, not on the smallest
colour mass or the number of colours.  All effective spaces in the
following definition are the first spaces
\(C_i^{\mathrm{old}}\), with total pin dimension at most \(K_1\).
They will not be enlarged.

Let \(\operatorname{pr}\) forget the channel columns and retain the
primal frames.  The marginal fact used below is exactly
\[
 \operatorname{pr}_*\pi=\operatorname{pr}_*\rawlaw.
\]
Thus a restriction of fixed positive pooled mass has endpoint law
at most a fixed multiple of \(\pi\), and its primal projection is
at most that same multiple of \(\operatorname{pr}_*\rawlaw\).
This does not assert a constant full-frame density relative to
\(\rawlaw\).

For a tag \(l\) and a diagonal bit \(q\), let
\(\Omega_{l,q,n}\) be the probability space of global plus and minus
vector slots on the star of \(l\), with the uniform law subject to
paired diagonal \(q\) in every component. Requiring its individual
vectors to be nonzero changes this law by exponentially small total
variation, which will be harmless.

\begin{definition}[Common prediction]
\label{def:prediction}
A prediction along a subsequence consists of common binary
functions
\[
 f_{i,n}(l,q,\cdot):\Omega_{l,q,n}\longrightarrow\F,
 \qquad i=1,2,
\]
a set of units of mass at least \(10^{-6}\), and, on every
unit in this set, choices
\[
 \lambda_i\in C_i^{\mathrm{old}},\qquad
 \alpha_i\in\F,\qquad
 {\cal E}_i\subseteq\F^b,\quad |{\cal E}_i|\le B_*,
\]
such that the following uniform statement holds.  There is
\(\varepsilon_n\to0\) for which, at both endpoints, at every
label outside \({\cal E}_i\), at every tag, and in every
permitted flavor, a uniform parameter atom \(x\) satisfies
\begin{equation}
 p_i(x)=T(\lambda_i,x)+\alpha_i a(x)
          +f_{i,n}(l,q,\operatorname{key}(x))
 \label{eq:prediction}
\end{equation}
except with probability at most \(\varepsilon_n\).
Here \(p_i=u_{i'}U_i\), where \(i'\) is the other endpoint.
The functions \(f_{i,n}\) are the same for all retained units.
The marks \(\lambda_i,\alpha_i,{\cal E}_i\) may depend on
the whole unit.
\end{definition}

The errors in \Cref{def:prediction} refer to uniform parameters
in each flavor, before conditioning on a tester bit.  Conditioning
on a possible fixed tester value only multiplies the error by a
fixed constant.  There are finitely many labels, tags and flavors,
with their numbers fixed independently of \(n\).

If a prediction exists, retain its marked population and
pigeonhole the values of the two \(\alpha_i\)'s, of
\[
 c=a(\lambda_1)+a(\lambda_2),
\]
and of the indicator \(\Delta=0\).
This loses a factor at most \(16\), leaving mass at least
\(6.25\cdot10^{-8}>10^{-12}\).
Pass to a subsequence on which every reference error probability
of every binary combination of \(f_1,f_2,q\) has a limit, for
each tag and bit.  This requires only finitely many comparisons.

For sequences of common functions, write \(f\simeq g\) if their
disagreement probability tends to zero on every
\(\Omega_{l,q,n}\), and identify sequences related in this way.
Then quotient their binary vector space by the span of the common
function with value \(q\).
Denote the resulting class by \([f]\), and put
\begin{equation}
 d_i=1+\alpha_i,\qquad e_i=(d_i,[f_i]).
 \label{eq:status-class}
\end{equation}
The pairs \(e_i\) are \emph{prediction classes}; a query's numerical
unary status records the scalar outcomes specified in
\Cref{sec:collision}.
The support of this pair of statuses is
\(\{i:e_i\ne0\}\).  These statuses are common to all retained
units, although the marked \(\lambda_i\)'s can vary.

\begin{lemma}[Exactification of empty status]
\label{lem:empty-status}
If \(e_1=e_2=0\), deleting \(o(1)\) unit mass makes
\begin{equation}
 p_i=T(\lambda_i,\cdot)+a
 \quad\text{on all of }\cX,\qquad i=1,2,
 \label{eq:exact-prediction}
\end{equation}
an exact identity on every retained unit.
The same conclusion holds on any marked population of fixed
positive mass satisfying \Cref{eq:prediction} with
\(\alpha_i=1\) and \(f_i\simeq\gamma_iq\); the threshold
\(10^{-6}\) in \Cref{def:prediction} is not needed here.
\end{lemma}

\begin{proof}
Empty status gives \(\alpha_i=1\) and
\(f_i\simeq\gamma_iq\) for a fixed \(\gamma_i\in\F\).
Restriction to any fixed positive population only changes the
primal marginal bound by a fixed factor, which is all the proof uses.
At any fixed tag, label and parameter point, the key depends
only on the primal frames.  Under \(\pi\) its law is therefore
the same as under \(\rawlaw\): the single-key reference up to
an exponentially small injectivity error, by
\Cref{lem:frame-images}.  The bounded \emph{primal} marginal
density transfers the reference disagreement of \(f_i\) and
\(\gamma_iq\) to an average \(o(1)\) parameter error on the
marked population.  No channel coordinate is tested in this
comparison.
There are only finitely many endpoint, tag and label choices.
Markov's inequality therefore allows deletion of \(o(1)\)
unit mass so that this error tends to zero uniformly on
the remaining units and choices.
Combining it with \Cref{eq:prediction} gives, at every
nonexceptional label,
\[
 p_i(x)+T(\lambda_i,x)+a(x)+\gamma_iq(x)=0
\]
with parameter error tending to zero in the generic flavor.

For fixed label and tag the left side is a Boolean polynomial
of base degree at most two.  A nonzero Boolean polynomial of
degree at most \(k\) has nonzero probability at least \(2^{-k}\)
under uniform bits.  For completeness, this follows by induction:
choose a coordinate with a nonzero derivative; the derivative
has degree at most \(k-1\), and whenever it is nonzero at least
one of the two values on that coordinate is nonzero.
The constant nonzero polynomial is the initial case.
Thus the displayed relation, once its error is below \(1/4\),
is exact for every allowed base value.

For any fixed allowed base value its dependence on the selector
has degree at most \(2j_*\).  It vanishes at every selector
outside a set of size at most \(B_*\).
The choice of \(b\) gives
\(B_*<2^{b-2j_*}\), so the same polynomial bound forces it
to vanish at the excluded labels also.

Use one shared-only base point with \(\eta_S=1\), the same at
every tag.  The sum of its \(g\) star atoms is zero in the cut
space.  Summing the exact relation over these atoms gives
\(0=\gamma_i g=\gamma_i\), because \(g\) is odd and
the first three terms are linear on \(\cX\).
Consequently \(\gamma_i=0\).
The point atoms span the cut space by its definition, proving
\Cref{eq:exact-prediction}.
\end{proof}

\subsection{A comparison for affine parameter slices}

The next elementary comparison is stated separately because the
global predictors in \Cref{def:prediction} need not be polynomial.
All tests in it are common functions of the displayed image
columns; they are not arbitrary orientation-dependent tests.

\begin{lemma}[Affine-slice comparison]
\label{lem:synthetic-slice}
Fix a selector, a flavor, and a fixed number \(d\) of affine
parameter variables.  Write a random affine coefficient map as
\[
 C(t,u)=p_s\otimes(t,z_0t+Lu),
 \qquad (t,u)\in\F\oplus\F^d.
 \]
Assume the flavor retains an \(n\)-bit \(Z\)-block absent from
\(E\).  Suppose we condition the affine coefficients on
\(C\trans EC=G\), where \(G\) is fixed, this event has probability
at least a fixed \(\beta>0\), and it depends on no \(Z\)
coefficient.
Let \({\cal T}_n\) be any common test of absolute value at most
one on the image column arrays \((P_eC,Q_eC)_e\), using a fixed
set of components.  Its empirical average over the conditioned
affine coefficients, with the primal frame sampled from
\(\operatorname{pr}_*\rawlaw=\operatorname{pr}_*\pi\), has variance
\[
 O(2^{-N/2}+2^{2d+1-n})
 \]
around its mean under uniform injective image column arrays
with internal Gram \(G\).
The constants may depend on the fixed column counts and
\(\beta\), but the estimate is uniform over the tests
\({\cal T}_n\).
\end{lemma}

\begin{proof}
The \(Z\)-coefficients remain independent uniform bits after
the conditioning.  The columns of \(C\) are independent
except with probability \(O(2^{d-n})\).
For two independent conditioned maps \(C,C'\), a relation
\[
 C(a,u)+C'(a',u')=0
 \]
has \(a=a'\), by its constant coordinate.  Its \(Z\)-part is
\[
 a(z_0+z_0')+L_Zu+L'_Zu'=0.
 \]
The matrix \([z_0+z_0',L_Z,L'_Z]\) is a uniform
\(n\times(2d+1)\) bit matrix.  A union over its nonzero
coefficient vectors bounds the failure of combined independence
by \((2^{2d+1}-1)2^{-n}\).
In particular, the two translation columns do not share a
forced nominal constant direction.

For fixed independent nominal columns, the restricted-image
transitivity in \Cref{lem:frame-images} gives the uniform
injective law with its specified Gram.  One batch consequently
has the claimed reference law, independently of the remaining
nominal coefficients.
For two batches, compared with two independent copies of this
reference law, the actual images have finitely many additional
mutual Gram entries prescribed.  Joint injectivity contributes
only \(O(2^{-N+O(d)})\) to the comparison: under independent
uniform vector slots, a dependence among a fixed number of
columns has this probability, and the fixed Gram conditioning
has bounded reciprocal probability.

Expand the mutual Gram indicator in binary characters, collapsing
any repeated identical entries.  A nonzero character has a
nonzero coefficient matrix on at least one interaction between
the two batches.  More explicitly the coefficient blocks pairing
first-batch plus vectors with second-batch minus vectors and
first-batch minus vectors with second-batch plus vectors give
cross-batch bit rank
\(N(\rank A+\rank B)\ge N\).
The internal Gram constraints, individual injectivity indicators,
and the two arbitrary tests are separate bounded weights on
the two batches; their normalization factors are fixed.
\Cref{lem:walsh} bounds every nontrivial character contribution
by \(O(2^{-N/2})\).  The same computation without the tests
gives the normalization of the mutual Gram conditioning.
There are only a fixed number of characters.  Thus the
conditional two-batch product expectation differs from the
product of the one-batch expectations by \(O(2^{-N/2})\),
uniformly over the nominal mutual Gram values.

Averaging the coefficients, and including the probability of
nominal dependence, proves the asserted second-moment estimate.
The one-batch mean differs from the reference mean only by
the already bounded nominal-dependence error.
\end{proof}

\begin{lemma}[Equal coefficients when the tensor difference vanishes]
\label{lem:zero-delta}
On a predicted branch with \(\Delta=0\) throughout,
\(\alpha_1=\alpha_2\).
\end{lemma}

\begin{proof}
Suppose otherwise.  By \Cref{lem:raw-intersections}, the two
canonical sparse representations \(w^{(i)}=(w_l^{(i)})_l\)
of the marked \(\lambda_i\)'s have the same value of \(\chi_*\).
Each has at most \(4D/g\) nonzero tag entries.
As \(\alpha_1\ne\alpha_2\), on every unit exactly one endpoint
has \(\chi_*(w^{(i)})+\alpha_i=1\).
Fix an endpoint \(i\) on a positive fraction of the units.
Pigeonhole the set
\[
 T_0=\{t:\eta_S(w_t^{(i)})=1\},
 \qquad |T_0|\le4D/g,
 \]
and choose a single selector outside the stored exceptions
on a positive fraction of this remaining population.
The latter is possible by averaging over selectors, since
\(B_*<2^{b-1}\).
All these restrictions have fixed positive mass in \(n\);
their endpoint primal marginal is at most
\(\Lambda_0\operatorname{pr}_*\rawlaw\) for a fixed \(\Lambda_0\).
No restriction made solely for this contradiction will be
used later to choose a channel-size threshold.

At each tag \(l\), use the pure flavor retaining exactly the
allowed \(O_{d,t}\) blocks with \(t\notin T_0\), setting all
\(S\) bits to zero and retaining the \(\#,Z\) blocks.
For a pure tag atom, \(a(x)=q(x)\) and every \(b_t(x)=0\).
The coefficient of its tester \(\eta_{d,t}\) in
\(T^0(\lambda_i,x)+\alpha_i a(x)\) is
\[
 a(\lambda_i)+b_t(\lambda_i)+\alpha_i
 =\chi_*(w^{(i)})+\eta_S(w_t^{(i)})+\alpha_i=1
 \]
on the retained \(t\)'s.  Hence \Cref{eq:prediction} gives
\begin{equation}
 f_i(l,q,\operatorname{key}(x))+q
 =p_i(x)+T^1(\lambda_i,x)
 \label{eq:pure-prediction}
\end{equation}
with uniform \(o(1)\) parameter error.

The right side of \Cref{eq:pure-prediction} is a sum of at
most
\begin{equation}
 s_1=(g-1)(h+2JK_1)
 \label{eq:synthetic-product-count}
\end{equation}
products of affine functions of the base variables.
Indeed \(p_i\) contributes at most \(h\) products in each
of the \(g-1\) star components.  Factoring the old
\(\lambda_i\)'s, whose total component rank is at most
\(K_1\), gives at most \(2J(g-1)K_1\) products for \(T^1\).
The selected pure tester retains at least
\begin{equation}
 \frac{g-1}{2}(g-4D/g)r_0>2s_1+1
 \label{eq:synthetic-tester-room}
\end{equation}
disjoint ordered bit pairs.  To check the last inequality,
use \(D=4000g\) and \(h=1000r_0\).  After division by
\((g-1)r_0\), the left side is \((g-16000)/2\) and the
right side is at most \(2000+4JK_1/r_0+o(1)\).
The fixed value \(g=10^9+1\) and the later choice of
sufficiently large \(r_0\) make the inequality hold.

Put \(m_1=2gs_1+1\), and introduce synthetic variables
\(\xi,\upsilon\in\F^{m_1}\) and
\[
 q_*=\sum_{j=1}^{m_1}\xi_j\upsilon_j.
 \]
For every component independently, choose uniform injective
maps of the affine coefficient columns
\((1,\xi,\upsilon)\) into the plus and minus ambient spaces
with full paired Gram
\[
 \pair{\text{plus}(1,\xi,\upsilon)}
       {\text{minus}(1,\xi',\upsilon')}
 =\sum_j\xi_j\upsilon'_j.
 \]
Such maps exist for large \(n\), since the number of columns
is fixed.  Their point evaluations form one synthetic array
of global component keys; all tag stars use the corresponding
entries of this same array.

We next prove two properties of this array with probability
tending to one.

\medskip\noindent
\emph{Shared-only parity.}
In a local shared-only test use the fixed selector and the
same point at every tag.  The sum of the tag atoms is zero
in \(\cX\).  Summing \Cref{eq:prediction} therefore gives
\[
 \sum_l f_i(l,q,\operatorname{key}_l)=0
 \]
except on \(o(1)\) parameter mass, uniformly on the retained
units.  Both tester bits have probability at least \(1/4\):
the bias of the sum of \(r_0\) independent bit-pair products
is \(2^{-r_0}\).

Apply \Cref{lem:synthetic-slice} with no affine variables
(\(d=0\)), the shared-only flavor, and internal Gram \((q)\).
The Gram event depends only on the fixed tester coordinates
and has probability at least \(1/4\).  The common test is
the displayed parity predicate of the global keys in all
components.  Its empirical frequency under the primal reference
law concentrates at its global reference probability.  The
endpoint primal marginal bound
\(\Lambda_0\operatorname{pr}_*\rawlaw\), and its frequency
\(1-o(1)\) on the retained population, force that reference
probability to tend to one.
At every point of the synthetic array the point-key marginal
is this reference, up to the exponentially small nonzero-vector
conditioning.  There are a fixed finite number of points.
Thus, with probability tending to one,
\begin{equation}
 \sum_l f_i(l,q_*,\operatorname{key}_l)=0
 \quad\text{at every synthetic point}.
 \label{eq:synthetic-parity}
\end{equation}

\medskip\noindent
\emph{Affine-product representations on every small slice.}
Fix sets \(I,J\subseteq\{1,\ldots,m_1\}\), each of size at
most \(2s_1+1\), and set all \(\xi\) and \(\upsilon\)
coordinates outside \(I,J\), respectively, to zero.
There are \(d=|I|+|J|\le4s_1+2\) remaining variables.
At a given tag, map this slice into its selected local pure
flavor by independent uniform affine functions for every
retained base bit, including independent translations.
Condition their full oriented Gram to equal the synthetic
slice Gram.

This conditioning has a fixed positive probability.
At the fixed selector \(E^\#\) vanishes on pairs.
There are at most \(2s_1+1\) active terms
\(\xi_j\upsilon'_j\) in the target Gram, so
\Cref{eq:synthetic-tester-room} lets us route them into
distinct retained ordered tester pairs, and set all other
tester affine coefficients to zero.  If \(t_0\) scalar
tester coordinates are prescribed, this assignment alone
has probability \(2^{-t_0(d+1)}>0\), independent of \(n\).
The \(Z\)-affine coefficients remain iid.
At every fixed slice point the unconditioned base input is
uniform; after conditioning, its density is bounded by the
fixed reciprocal Gram probability.

Consequently \Cref{eq:pure-prediction} holds simultaneously
at all points of the slice except on \(o(1)\) of these
parametrizations, uniformly on the marked population.
Its right side pulls back to a sum of at most \(s_1\)
affine products.  Consider only the event
\begin{quote}
the array of values of the common function \(f_i+q_*\) on
this specified slice admits such an affine-product
representation.
\end{quote}
This is a common test of the image column arrays.  It contains
neither the mark \(\lambda_i\) nor any coefficients of the
unit-specific representation.  Its empirical success
probability is \(1-o(1)\) on the retained population.
\Cref{lem:synthetic-slice} and the bounded primal marginal
therefore force its probability under the synthetic slice
law to tend to one.  That law is also the marginal of the
corresponding slice in the full synthetic array, by restricted
frame transitivity.

There are only finitely many choices of tag and slice, with
the number fixed independently of \(n\).  We may therefore
require all these events and \Cref{eq:synthetic-parity}
simultaneously.  Notice that we embedded only the small
slices locally; no local realization of the entire synthetic
tester was required.

Choose one synthetic array on which all these properties hold,
and define
\[
 F_l(\xi,\upsilon)
   =f_i(l,q_*,\operatorname{key}_l)+q_*.
 \]
For each tag define its Boolean cross-coefficient matrix by
\[
 C_l[j,k]=F_l(e_j,e_k)+F_l(e_j,0)
                 +F_l(0,e_k)+F_l(0,0).
 \]
On a coordinate slice where \(F_l\) is a sum of \(s_1\)
affine products, this is the matrix of \(\xi_j\upsilon_k\)
coefficients.  One affine product contributes a sum of two
rank-one matrices to this cross matrix.  Every
\((2s_1+1)\)-square minor of \(C_l\) is therefore singular,
so \(\rank C_l\le2s_1\).
On the other hand \Cref{eq:synthetic-parity}, together with
odd \(g\), gives \(\sum_lF_l=q_*\) and hence
\[
 \sum_l C_l=I_{m_1}.
 \]
Rank subadditivity would imply
\(m_1\le\sum_l\rank C_l\le2gs_1\), contradicting
\(m_1=2gs_1+1\).  This proves \(\alpha_1=\alpha_2\).
\end{proof}

\subsection{Compatible tables and the prepared alternatives}

For an admissible small table, a contraction on a marked
\(\lambda_i\) means its table contraction.  This is well defined
because \(\lambda_i\in C_i^{\mathrm{old}}\) and all first effective
spaces are retained.  We use the following phase requirements:
\begin{equation}
\begin{array}{ll}
\displaystyle
u_{B,S^\circ}(\Delta_A)+u_{A,S^\circ}(\Delta_B)
       =|S^\circ|d,
&
\begin{gathered}
\text{if the nonzero }e_i\text{ all coincide},\\[-1mm]
S^\circ=\{i:e_i\ne0\}\ne\varnothing,\quad
 d=\text{their first bit};
\end{gathered}
\\[4mm]
u_{B,j}(\Delta_A)=c,\quad u_{A,i}(\Delta_B)=c
\quad(i,j=1,2),
&\text{if }e_1=e_2=0.
\end{array}
\label{eq:phase-compatibility}
\end{equation}
In all other status cases there is no additional phase condition.
The four equations in the second line all use the \emph{same}
previously fixed bit \(c=a(\lambda_1)+a(\lambda_2)\).
For a candidate table these contractions are computed from its
entries and the numerical coordinates of the marks.

\begin{proposition}[Status preparation]
\label{prop:status-preparation}
There are constants \(\delta_*,c_*>0\), fixed before the selector
and channel sizes, with the following property.  Along a subsequence,
the original pooled unit law has a restriction of mass at least
\(\delta_*\) whose first pin dimension is at most \(K_1\) and whose
normalized leaves satisfy \Cref{eq:final-leaf-minentropy}.  One of
the following descriptions applies.
\begin{enumerate}[label=\textup{(\roman*)}]
\item There is no prediction as in \Cref{def:prediction} on the
first pooled law before the final negligible leaf trimming.  For
independent draws from the prepared law, the probability of an
admissible injecting actual small table is greater than \(.99\).
\item A prediction with fixed \(\alpha_i,c\) and common classes
\(e_i\) is retained.  Its marks and exceptions are preserved, and
\Cref{eq:prediction} has uniform \(o(1)\) error on every retained
unit.  If both classes vanish, \Cref{eq:exact-prediction} is exact
on every unit.  For independent draws from the prepared law,
the probability of an admissible injecting actual table satisfying
\Cref{eq:phase-compatibility} is at least \(1/64\).
\end{enumerate}
In particular the corresponding successful pair event, measured
in two draws from the original pooled law, has mass at least
\(c_*\).  One may take
\begin{equation}
 \delta_* = \frac{10^{-6}\delta_{\mathrm{ph}}}{64},
 \qquad c_* = \frac{\delta_*^2}{64},
 \qquad 0<\eta_{\mathrm{col}}<\frac{c_*}{4},
 \label{eq:preparation-early-margins}
\end{equation}
where \(0<\delta_{\mathrm{ph}}\le1\) is the early retained-mass
constant in \Cref{thm:phase-alternative}.  The colour bound in
\Cref{eq:preparation-early-margins} can be chosen using the initial
\(M\) and early geometric constants only; the marginal constant
\(L\) of the second distribution test does not enter it.

After forgetting colour and other marks, the prepared law
\(\sigma'\) satisfies
\begin{equation}
 \sigma'_i\le\frac{M}{\delta_*}\pi\quad(i=1,2),
 \qquad
 \sigma'\le 2^{(D+o(1))N}\rawlaw^2.
 \label{eq:prepared-marginals}
\end{equation}
The endpoint primal marginals consequently have the same constant
bound relative to \(\operatorname{pr}_*\rawlaw\).  Every normalized
leaf has joint density at most \(2^{O(N)}\rawlaw^2\).  The successful
table event is detected by a finite list of numerical tables and
separate unary orientation flags, together with leaf-pair tests.
The size of this list is fixed in \(n\), but need not be early.
\end{proposition}

\begin{proof}
Consider first the pooled law after the first peeling and
intersection deletion.  If no prediction exists on any subsequence,
apply \Cref{lem:injection} to independent draws with no additional
acceptance condition.  Its failure tolerance can be chosen below
\(1/128\), so the success probability is greater than \(.99\) for
large \(n\).  The actual table is automatically admissible.  Later
leaf trimming removes only \(o(1)\) mass.  A prediction on the
trimmed law with mass bounded strictly above \(10^{-6}\) would give
one on the preceding law, by multiplying its mass by the retained
proportion and preserving its common functions and unitwise error.
This is the form of the no-prediction conclusion used below.

Otherwise retain a prediction population and pigeonhole
\(\alpha_1,\alpha_2,c\) and the indicator \(\Delta=0\), as above.
Its relative mass is at least \(10^{-6}/16\).  If the statuses
are empty, perform the \(o(1)\) deletion in
\Cref{lem:empty-status}.  All subsequent restrictions then preserve
the exact identities \Cref{eq:exact-prediction}.

Suppose \(\Delta=0\).  By \Cref{lem:zero-delta},
\(\alpha_1=\alpha_2\).  If precisely one status is active, its
first bit cannot be \(1\): that would give \(\alpha_i=0\) there
and \(\alpha_{i'}=1\) at the inactive endpoint.  Thus the right
side of the first line of \Cref{eq:phase-compatibility} is zero.
With two equal active statuses it is also zero, because \(2d=0\)
over \(\F\).  Every left side vanishes when \(\Delta=0\), so
all these conditions hold automatically.

If both statuses are empty and \(c=1\), then
\(U_1\lambda_1=U_2\lambda_2\), while the two identities
\Cref{eq:exact-prediction} give the affine-gradient equations
for an internal hole and \(c=1\) gives its unequal \(a\)-values.
This contradicts the assumption that the pair is a unit.
Consequently \(c=0\), and all four required phase bits vanish.
Thus compatibility is automatic also in this case.

Suppose next that \(\Delta\ne0\).  Apply
\Cref{thm:phase-alternative} to the pooled marked law.  It retains
at least the early fraction \(\delta_{\mathrm{ph}}\), and gives
either full mixing or the identity
\[
 F(A,B)=u_{B,\{1,2\}}(\Delta_A)=0
 \quad\text{for every two retained units }A,B,
\]
together with exponentially small singleton character means.
For clarity, write the four phase bits as
\[
 X_j=u_{B,j}(\Delta_A),\qquad
 Y_i=u_{A,i}(\Delta_B),\qquad i,j\in\{1,2\}.
\]
Full mixing bounds the absolute mean of every nonconstant
character of \((X_1,X_2,Y_1,Y_2)\) by \(.005\).  Binary Fourier
inversion therefore gives every prescribed four-bit pattern
probability at least
\[
 \frac{1-15(.005)}{16}>\frac1{32}.
\]
Each condition in \Cref{eq:phase-compatibility} admits such a
pattern.  In particular the empty-status pattern is exactly
\((c,c,c,c)\).

In the other alternative, \(X_1=X_2\) and \(Y_1=Y_2\)
identically.  If there is one active status, the required bit is
\(X_i+Y_i=d\); its nontrivial character has singleton support in
both units, so its probability is \(1/2+o(1)\).  If two active
statuses coincide, the condition is
\(F(A,B)+F(B,A)=0\), which holds identically.  With empty statuses
put \(X=X_1=X_2\) and \(Y=Y_1=Y_2\).  The characters of
\(X,Y,X+Y\) all have singleton representatives.  Hence
\begin{align*}
 \PP\{X_1=X_2=Y_1=Y_2=c\}
 &=\frac14\left(1+(-1)^c\E(-1)^X
                    +(-1)^c\E(-1)^Y+\E(-1)^{X+Y}\right)\\
 &=\frac14+o(1)>\frac1{32}.
\end{align*}
This explicitly requires the same fixed \(c\) on both sides.
The remaining status cases impose no condition.  Thus all branches
have conditional compatibility probability at least \(1/32\) for
large \(n\), in agreement with \Cref{tp:phase-compatibility}.

Apply the unconditional form of \Cref{lem:injection} to this
retained law, with failure below \(1/128\).  The probability of a
compatible injecting actual table is then at least
\[
 \frac1{32}-\frac1{128}-o(1)>\frac1{64}.
\]
This is subtraction of exceptional mass; it makes no independence
assumption about injection and compatibility.  The injection
estimate uses the endpoint bound relative to \(\pi\), together
with the unchanged primal marginal and the allowed thinning loss
in the full-frame estimates.  Its tolerance is fixed before \(h\).

We now justify the claimed leaf estimates without adding any pins.
Let \(w_\lambda\) be the original actual weights of the first
leaves, and let \(q_\lambda\) be the fraction of leaf \(\lambda\)
retained by all restrictions just made.  Discard the retained parts
with \(q_\lambda<2^{-\zeta N}\).  Their total unnormalized mass
is at most
\[
 \sum_{\lambda:q_\lambda<2^{-\zeta N}}
       w_\lambda q_\lambda\le 2^{-\zeta N}.
\]
On each remaining leaf, division by \(q_\lambda\) weakens
\Cref{eq:leaf-minentropy} to
\[
 2^{\zeta N}2^{-(1-\zeta)tN}
 \le 2^{-(1-2\zeta)tN}\qquad(t\ge1),
\]
as required.  Its joint density gains at most \(2^{\zeta N}\)
and remains \(2^{O(N)}\rawlaw^2\).  The pin spaces and all old
marks are unchanged.  This exponentially small deletion changes
the successful pair probability by \(o(1)\), already allowed above.

The first peeling loses \(o(1)\) mass.  On a predicted branch,
the finite initial selection costs at most \(16\cdot10^6\), and
the phase restriction costs at most \(1/\delta_{\mathrm{ph}}\).
All further losses tend to zero; for large \(n\) their combined
retained factor, including first peeling, is at least \(1/4\).
This proves the unit-mass lower bound \(\delta_*\) in
\Cref{eq:preparation-early-margins}.  The no-prediction branch
retains \(1-o(1)\) mass and satisfies the same lower bound.
Multiplying the conditional pair-success bound by
\(\delta_*^2\) proves the original-scale bound \(c_*\).
The density statements follow by dividing the initial bounds by
the retained unit mass and using \Cref{eq:thinning-density}.
All marginal constants required by the phase and injection steps
are fixed in advance, for example below \(2M/\delta_*\), so their
channel-size thresholds can be imposed after these early margins
have been selected.

Finally, every numerical table has bounded coordinate dimensions.
The marked tensors are recorded by their coordinates in the first
effective bases, and the exception sets have bounded descriptions.
For a fixed leaf pair and a fixed numerical option, the two
admissibility conditions are separate unary orientation flags,
since the pin images are fixed on each leaf.  Fixing numerical
mark coordinates is also a unary flag.  The phase requirements
then depend only on the table entries and those coordinates.
An actual successful table supplies at least one successful
option in this finite list.  This proves the asserted finite-data
form without selecting one table during preparation or conditioning
the two draws jointly on table success.
\end{proof}

\begin{corollary}[Removing equal colours]
\label{tp:preparation-colours}
Suppose the original colour masses are \(p_\gamma\le
\eta_{\mathrm{col}}\).  The successful prepared pair event of
\Cref{prop:status-preparation}, with the further requirement that
its two colours differ, has mass at least
\(c_*-\eta_{\mathrm{col}}>3c_*/4\) in the original pooled law.
\end{corollary}

\begin{proof}
Two independent original draws have equal colours with probability
\[
 \sum_\gamma p_\gamma^2\le
 \eta_{\mathrm{col}}\sum_\gamma p_\gamma=\eta_{\mathrm{col}}.
\]
Subtract this upper bound from the original-scale successful mass
\(c_*\).  Equivalently, in a restriction of mass \(\delta\), the
equal-colour probability is at most
\(\eta_{\mathrm{col}}/\delta^2\).  This calculation applies also
to randomized colours, which are simply part of the original
marked probability space.  It uses actual colour probabilities
and imposes no lower bound on them.
\end{proof}

\begin{remark}[No feedback from later accuracies]
\label{rem:preparation-accuracies}
The margins \(\delta_*,c_*\), the choice of
\(\eta_{\mathrm{col}}\), and the injection tolerance are fixed
before \(h\).  At this stage success means that some numerical
table passes.  The later choice of one option from a finite list
may have a much smaller mass depending on \(h\), which is
irrelevant to the earlier injection estimate.  The errors in the
affine-slice comparison and in subsequent fixed-test limits are
handled by increasing \(n\); their slice counts, partitions and
flag thresholds do not change the channel size or pin budget.
Only the first peeling has been used throughout.
\end{remark}
